\chapter{Conclusion and Outlook}\label{sec:conclusion}

This thesis extends MetaConfigurator with an integrated \gls{rdf} authoring workflow that connects schema-driven \gls{json} editing with Semantic Web technologies. The resulting workflow combines four capabilities in one interface: (1) transformation from \gls{json} to \gls{jsonld} via \gls{rml}, (2) semantic editing through context and triple views, (3) in-browser \gls{sparql} querying with validation support, and (4) interactive \gls{kg} visualization with editing actions. \gls{ai}-assisted generation of \gls{rml} mappings and \gls{sparql} queries further lowers the entry barrier for non-expert users, while keeping humans in control through explicit review before execution.

The application chapter demonstrates this workflow on \gls{mof} synthesis data. Laboratory protocol records are transformed into ontology-grounded \gls{rdf} resources (\gls{chmo}, \gls{chebi}, and \gls{obi}), then queried and explored as a graph. In practice, this shows that semantic data authoring can be carried out in one environment, reducing tool switching and improving interoperability.

\section{Current Limitations}

Despite these contributions, several limitations remain.

\begin{enumerate}
  \item \textbf{Text-view serialization pattern after the first RDF edit:} After the first data edit in the \gls{rdf} panel, the JSON-LD rendered in the Text View is regenerated from the internal \gls{rdf} statements using a predefined serialization pattern. The resulting text can differ from the originally authored document even when the represented semantics stay equivalent. As shown in Listing~\ref{lst:conclusion_jsonld_example}, removing the triple:
  
  \begin{quote}
  \small
  \texttt{(local:variable\_element\_size , \texttt{schema:name} , "element\_size")}
  \end{quote}
    
  leads to additional textual changes such as renaming the \texttt{schema} prefix to \texttt{sch} and reordering literal fields (e.g., \texttt{@value} and \texttt{@type}). This behavior makes it harder to preserve original authoring style and document-level traceability.

  \item \textbf{Remaining error risk in AI-assisted artifacts:} \gls{ai}-generated mappings and queries can still contain semantic or structural errors. User review therefore remains essential, especially for domain-critical data and ontology alignment decisions.

  \item \textbf{Scalability limits for browser-side workflows:} Scalability is still bounded by browser-side processing constraints for larger graphs, especially during interactive visualization and query-result rendering. To mitigate this, two configurable limits are already available:
  \begin{itemize}
    \item \texttt{maximumTriplesToShow}: limits how many triples are rendered in the \gls{rdf} panel to keep table rendering and interaction responsive.
    \item \texttt{maximumNodesToVisualize}: limits how many nodes are rendered in the \texttt{Visualization} dialog to avoid unstable layout computation and slow graph interactions.
  \end{itemize}
  Both settings can be adjusted according to dataset size and client performance.

  \item \textbf{Ontology and IRI selection still requires expert knowledge:} Even with integrated mapping support, high-quality transformation from \gls{json} to \gls{rdf} still depends on choosing suitable ontologies, classes, predicates, and identifier patterns. Producing semantically robust mappings often requires domain and Semantic Web expertise (e.g., selecting appropriate ontology terms and stable \glspl{iri}), which remains a substantial barrier for non-expert users.
\end{enumerate}

\begin{lstlisting}[float=htbp,style=json, caption={De-serilaized JSON-LD in Listing~\ref{lst:background_json_example} after the first edit in the \gls{rdf} Panel results in renaming \texttt{schema} prefix and reordering \texttt{@value} and \texttt{@type} fields}, label={lst:conclusion_jsonld_example}]
{
  "@context": {
    "m4i": "http://w3id.org/nfdi4ing/metadata4ing#",
    "unit": "http://qudt.org/vocab/unit/",
    "local": "https://www.example.org/",
    "sch": "https://schema.org/"
  },
  "@graph": [
    {
      "@id": "local:run_fenics_simulation",
      "@type": "m4i:Method",
      "m4i:hasParameter": {
        "@id": "local:variable_element_size"
      },
      "m4i:investigates": {
        "@id": "local:variable_max_von_mises_stress"
      }
    },
    {
      "@id": "local:variable_element_size",
      "@type": "sch:PropertyValue",
      "sch:unitCode": {
        "@id": "unit:M"
      },
      "sch:value": {
        "@value": "2.5E-2",
        "@type": "http://www.w3.org/2001/XMLSchema#double"
      }
    },
    {
      "@id": "local:variable_max_von_mises_stress",
      "@type": "sch:PropertyValue",
      "sch:name": "max_von_mises_stress",
      "sch:unitCode": {
        "@id": "unit:MegaPA"
      },
      "sch:value": {
        "@value": "2.997915075586333E8",
        "@type": "http://www.w3.org/2001/XMLSchema#double"
      }
    }
  ]
}
\end{lstlisting}

\section{Outlook}

Future work should focus on five directions:

\begin{enumerate}
  \item \textbf{Preserve textual form during synchronization:} Future work should improve synchronization using path-aware or diff-based updates from statement changes back to the Text View. This is a syntactic and formatting concern: the semantic content is already preserved correctly, but the rendered textual form can drift from the originally authored style.

  \item \textbf{Integrate semantic quality checks:} Semantic quality assurance should be strengthened through integrated \gls{shacl} validation in the authoring loop, so users can check constraints before export and publication.

  \item \textbf{Improve scalability mechanisms:} Progressive graph loading, query-result size control, and graph summarization should be added to better support larger datasets.

  \item \textbf{Support ontology discovery from current data context:} An additional extension would be an integrated recommendation feature that suggests suitable ontologies, classes, and predicates based on the currently loaded \gls{json} data. This could reduce ontology-selection effort during mapping and help non-expert users choose semantically appropriate vocabularies.

\end{enumerate}

Overall, the presented extension shows that user-centered semantic authoring in scientific workflows can be carried out within a single application, from structured \gls{json} input to ontology-aware querying and \gls{kg} analysis.
